\]
y el máximo valor absoluto del momento interno es:
\[
|M|_{\max}=\boxed{\qty{20}{\kilo\newton\metre}}.
\]

\Needspace{5\baselineskip}\section*{Problema 3}
Viga simplemente apoyada de $L=\qty{3}{\metre}$ sometida a carga triangular $q_0=\qty{10}{\kilo\newton\per\metre}$ máxima en $A$ y nula en $B$.

\subsection*{Resultado}
\[
Q=\qty{15}{\kilo\newton},\qquad \bar x=\qty{1}{\metre},
\]
\[
\boxed{R_A=\qty{10}{\kilo\newton},\qquad R_B=\qty{5}{\kilo\newton}}.
\]

\Needspace{5\baselineskip}\section*{Problema 4}
Barra de acero:
\[
N=\qty{120}{\kilo\newton},\quad A=\qty{800}{\milli\metre\squared},
\]
\[
L=\qty{2}{\metre},\qquad E=\qty{210}{\giga\pascal}.
\]

\subsection*{Resultado}
\[
\boxed{\sigma=\qty{150}{\mega\pascal}}
\]
\[
\boxed{\varepsilon=7.143\times10^{-4}}
\]
\[
\boxed{\Delta L=\qty{1.43}{\milli\metre}}.
\]

\Needspace{5\baselineskip}\section*{Problema 5}
Viga simplemente apoyada de $L=\qty{4}{\metre}$ sometida únicamente a un momento exterior de $M_0=\qty{12}{\kilo\newton\metre}$ aplicado en el centro.

Calcular las reacciones.

\subsection*{Resultado}
En magnitud:
\[
\boxed{|R_A|=|R_B|=\qty{3}{\kilo\newton}}
\]
con sentidos opuestos. El momento aplicado genera un salto en el diagrama de $M$.

\Needspace{5\baselineskip}\section*{Problema 6}
Se conoce:
\[
V(x)=5-2x^2\quad[\mathrm{kN}].
\]
Determinar $q(x)$.

\subsection*{Resultado}
\[
q(x)=-\frac{\dd V}{\dd x}=\boxed{4x\quad[\mathrm{kN/m}]}.
\]

% ============================================================
